\par\noindent\textbf{(1)} 平衡时 \(n_{A_2}=n_0(1-\alpha)\)，\(n_A=2n_0\alpha\)，总物质的量 \(n_{\mathrm{tot}}=n_0(1+\alpha)\)。因此摩尔分数 \(y_A=2\alpha/(1+\alpha)\)、\(y_{A_2}=(1-\alpha)/(1+\alpha)\)，总压强因解离膨胀。
\par\noindent\textbf{(2)} 理想气体化学势满足 \(\mu_i=\mu_i^\circ+RT\ln(y_iP/P^\circ)\)。平衡条件 \(2\mu_A-\mu_{A_2}=0\) 给 \(K_p=(p_A/P^\circ)^2/(p_{A_2}/P^\circ)=4\alpha^2(P/P^\circ)/(1-\alpha^2)\)。必须保留标准压强使 \(K_p\) 无量纲。
\par\noindent\textbf{(3)} 由 \(K_p(1-\alpha^2)=4\alpha^2(P/P^\circ)\)，得到 \(\alpha^2=K_p/[K_p+4P/P^\circ]\)。在定温 \(K_p>0\) 下，\(P\) 增大使解离度降低，符合反应气体分子数从1增至2的Le Chatelier趋势。
\par\noindent\textbf{(4)} 当 \(P/P^\circ\ll K_p\) 时 \(\alpha\simeq1-2P/(K_pP^\circ)\)，几乎完全解离；高压 \(P/P^\circ\gg K_p\) 时 \(\alpha\simeq\tfrac12\sqrt{K_pP^\circ/P}\)。如果活度系数显著偏离1，真实高压气体需要逸度而不是简单分压。
